\section{Training Procedure}

\begin{frame}[label=training_cycle]{The Training Cycle}
   \centering
    
    % Wraps the TikZ diagram to ensure it fits the slide width perfectly
    \hspace*{-5mm}
    \begin{tikzpicture}[
        box/.style={rectangle, draw=blue!60!black, fill=blue!10, rounded corners, thick, minimum width=2.8cm, minimum height=1cm, align=center, font=\scriptsize\bfseries},
        final/.style={rectangle, draw=blue!60!black, fill=blue!10, rounded corners, thick, minimum width=1.8cm, minimum height=1cm, align=center, font=\scriptsize\bfseries},
        decision/.style={diamond, draw=blue!60!black, fill=blue!10, thick, align=center, aspect=1.4, font=\scriptsize\bfseries, inner sep=1pt},
        arrow/.style={-Latex, thick, blue!70!black, rounded corners=4pt}
    ]
        % Main horizontal axis (Y = 0)
        \node[box] (init) at (-0.5, 0) {Initialization \\ (Point Cloud)};
        \node[box] (select) at (3, 0) {1. Select Image \\ \& Camera Pose};
        \node[decision] (loss) at (8, 0) {3. Loss Function \\ Converged?};
        \node[final, fill=green!10, draw=green!50!black] (end) at (11.5, 0) {Final \\ 3D Scene};
        
        % Top and Bottom nodes forming the loop
        \node[box] (forward) at (5.5, 2.5) {2. Forward Pass \\ (Render Image)};
        \node[box] (backward) at (5.5, -2.5) {4. Backward Pass \\ (Update Attributes)};
        
        % Horizontal entry and exit arrows
        \draw[arrow] (-2.3, 0) -- (init.west);
        \draw[arrow] (init.east) -- (select.west);
        \draw[arrow] (loss.east) -- node[above, font=\tiny, text=blue!80!black] {Yes} (end.west);
        
        % Arrows forming the rectangular loop (|- and -| create the 90-degree corners)        
        \draw[arrow] (select.north) |- (forward.west);
        \draw[arrow] (forward.east) -| (loss.north);
        \draw[arrow] (loss.south) |- node[pos=0.25, right, font=\tiny, text=blue!80!black] {No} (backward.east);
        \draw[arrow] (backward.west) -| (select.south);
        
    \end{tikzpicture}%
    \pdfcomment[opacity=0]{
        - Let's examine point cloud initialization in detail
    }
\end{frame}

\subsection{Point Cloud Initialization}

\begin{frame}{Point Cloud Initialization}
    \begin{columns}[T]
        % Left Column: Spotlighted Concepts (0.35 split)
        \begin{column}{0.35\textwidth}
            \textbf{Initialization Steps}
            \vspace{1em}
            
            \alt<1>{\textbf{\textcolor{blue}{Structure from Motion}}}{\textcolor{gray}{Structure from Motion}}\par\vspace{1em}
            \alt<2>{\textbf{\textcolor{blue}{Camera Matrix}}}{\textcolor{gray}{Camera Matrix}}\par\vspace{1em}
            \alt<3>{\textbf{\textcolor{blue}{Primitive Initialization}}}{\textcolor{gray}{Primitive Initialization}}
        \end{column}

        % Right Column: Synchronized Explanations (0.6 split)
        \begin{column}{0.6\textwidth}
            \textbf{Details}
            \vspace{1em}
            
            \begin{overlayarea}{\textwidth}{3cm}
                \justifying 
                \only<1>{
                    The 3DGS algorithm requires initial geometry, which is commonly derived from a set of input images.\\
                    \vspace{1em}
                    Because camera poses are usually unknown, Structure from Motion (SfM) software like COLMAP is used to solve a geometric optimization problem, providing an accurate 3D point cloud.
                }
                \only<2>{
                    As another result from SfM, each input image is assigned a $3\times3$ camera matrix~$W_{rot}$ which contains the camera's orientation in world space.
                }
                \only<3>{
                    The SfM points act as the origins for the initial 3D Gaussians. 
                    \begin{itemize}
                        \justifying
                        \item They are initialized as perfect spheres.
                        \item Radii are deduced from their three nearest neighbors.
                        \item Opacity is set to $0.1$ and color to any RGB value.
                    \end{itemize}
                }
            \end{overlayarea}
        \end{column}
    \end{columns}
    \pdfcomment[opacity=0]{
        - Usually, in 3D Reconstruction, we are faced with an initial problem of not knowing where in world space the input images were taken. \textCR
        - Introduce world space, view space, image space \textCR 
        \textCR
        - After the initialization phase we progress to the forward pass.
    }
\end{frame}


\subsection{Forward Pass}

\begin{frame}[label=forwardpass_pipeline]{Forward Pass: Pipeline Architecture}
    \centering
    \vspace{1.5em}
    \hspace*{-7mm}
    \resizebox{1.1\linewidth}{!}{%
    \begin{tikzpicture}[
        phase/.style={rectangle, draw=blue!80!black, fill=blue!5, thick, rounded corners=4pt, minimum width=2.4cm, minimum height=1.2cm, align=center, font=\scriptsize\bfseries},
        arrow/.style={-Latex, thick, blue!60!black}
    ]
        % --- Phase Nodes ---
        \node[phase] (p1) at (0,0) {Phase 1\\Attributes};
        \node[phase] (p2) at (2.8,0) {Phase 2\\3D Matrix};
        \node[phase] (p3) at (5.6,0) {Phase 3\\Projection};
        \node[phase] (p4) at (8.4,0) {Phase 4\\Tile Mapping};
        \node[phase] (p5) at (11.2,0) {Phase 5\\Radix Sort};
        \node[phase] (p6) at (14.0,0) {Phase 6\\Rasterization};
        
        % --- Arrows ---
        \draw[arrow] (p1) -- (p2);
        \draw[arrow] (p2) -- (p3);
        \draw[arrow] (p3) -- (p4);
        \draw[arrow] (p4) -- (p5);
        \draw[arrow] (p5) -- (p6);
        
        % --- Top Brackets: Spatial Transformations ---
        \draw[gray!80!black, thick] (-1.2, 0.7) -- (-1.2, 0.9) -- (4.0, 0.9) -- (4.0, 0.7);
        \node[above, font=\small\bfseries] at (1.4, 0.95) {3D World Space};
        
        \draw[gray!80!black, thick] (4.4, 0.7) -- (4.4, 0.9) -- (6.8, 0.9) -- (6.8, 0.7);
        \node[above, font=\small\bfseries] at (5.6, 0.95) {World $\rightarrow$ View $\rightarrow$ Image};
        
        \draw[gray!80!black, thick] (7.2, 0.7) -- (7.2, 0.9) -- (15.2, 0.9) -- (15.2, 0.7);
        \node[above, font=\small\bfseries] at (11.2, 0.95) {2D Image Space};
        
        % --- Bottom Brackets: GPU Architecture Modes ---
        \draw[gray!80!black, thick] (-1.2, -0.7) -- (-1.2, -0.9) -- (9.6, -0.9) -- (9.6, -0.7);
        \node[below, font=\small\bfseries] at (4.2, -0.95) {One GPU Thread per 3D/2D Gaussian};
        
        \draw[gray!80!black, thick] (10.0, -0.7) -- (10.0, -0.9) -- (12.4, -0.9) -- (12.4, -0.7);
        \node[below, font=\small\bfseries, align=center] at (11.2, -0.95) {Global Sync\\(Sorting)};
        
        \draw[gray!80!black, thick] (12.8, -0.7) -- (12.8, -0.9) -- (15.2, -0.9) -- (15.2, -0.7);
        \node[below, font=\small\bfseries, align=center] at (14.0, -0.95) {One GPU Thread\\per Pixel};
        
    \end{tikzpicture}%
    }
    \pdfcomment[opacity=0]{
        - We can subdivide the forward pass into 6 phases. \textCR
        - What's also important from now on is looking at which space a phase operates in, which you can see above the phases as well as which mode the GPU is running in, which you can see below the phases. \textCR
        - We will look at each phase in detail on the following slides, but the general path is: we project 3D Gaussian to 2D Gaussian, sort them and finally as the rasterization step, splat them against the camera sensor.
    }
\end{frame}

%\begin{frame}{Phase 1: Preparing Attributes}
\begin{frame}{Forward Pass: Phases 1 \& 2}

    \textbf{Phase 1: Preparing Constituent Attributes}
    \begin{itemize}
        \justifying
        \item The GPU is running one GPU thread per 3D Gaussian.
        \item Raw values optimized during training are converted from latent space back into physical space.
        \item Opacity is bounded to $(0,1)$ using a sigmoid function, scale uses the exponential function, and the rotational quaternion is normalized.
    \end{itemize}
    
    \vspace{1em}
    
    \pause
    \textbf{Phase 2: Reconstructing the 3D Covariance Matrix}
    \begin{itemize}
        \justifying
        \item The algorithm mathematically rebuilds the 3D geometric shape.
        \item It combines the scaling matrix $S$ and rotational matrix $R$ to form the full 3D covariance matrix:
        $$ \darkmathhighlight{\Sigma_{3D} = RSS^{T}R^{T}} $$
    \end{itemize}
\end{frame}

\begin{frame}{Forward Pass: Phase 3}
    % Overlay 1: Text only
    \only<1>{
        \textbf{Phase 3: Projection from 3D to 2D}
        \begin{itemize}
            \justifying
            \item Converts 3D Gaussians (ellipsoids) in world space into flat 2D Gaussians (ellipses) in image space.
            \item Employs a linear affine approximation (the Jacobian matrix $J$) of a standard non-linear pinhole projection.
            \item The 3D matrix is rotated into view space via $W_{rot}$ and then projected:
            $$ \darkmathhighlight{\Sigma_{2D} = J W_{rot} \Sigma_{3D} W_{rot}^{T} J^{T}} $$
        \end{itemize}
    }
    
    % Overlay 2: Graphic only (replaces the text on the same frame number)
    \only<2>{
        \begin{figure}
            \centering
            \begin{tikzpicture}[scale=0.9, every node/.style={transform shape}]
        
                % ==========================================
                % 1. 3D GAUSSIAN (WORLD/VIEW SPACE)
                % ==========================================
                \coordinate (mu3D) at (0, 0);

                % Title
                \node[blue!80!black, font=\Large\bfseries, align=center] at (0, 3) {3D World Space};

                % 3D Fuzzy Blob
                \foreach \i in {1,2,3,4,5,6} {
                    \pgfmathsetmacro{\opac}{0.85 - \i*0.13}
                    \fill[green!50!black, opacity=\opac, rotate around={20:(mu3D)}]
                        (mu3D) ellipse ({2.0 - \i*0.3} and {1.3 - \i*0.2});
                }

                % 3D Coordinate Stub representing Covariance (\Sigma_{3D})
                % Arrows point along the principal axes of the 3D ellipsoid
                \draw[-Latex, very thick, blue!80!black] (mu3D) -- ++(20:1.8) node[right, font=\scriptsize] {}; % Major axis X
                \draw[-Latex, very thick, blue!80!black] (mu3D) -- ++(110:1.1) node[above, font=\scriptsize] {}; % Minor axis Y
                \draw[-Latex, very thick, blue!80!black] (mu3D) -- ++(230:0.9) node[below left, font=\scriptsize] {}; % Depth axis Z

                % Center Point and Label
                \node[circle, fill=black, inner sep=1.8pt] at (mu3D) {};
                \node[below right, font=\small\bfseries, align=left, fill=white, fill opacity=0.7, text opacity=1, inner sep=2pt, rounded corners] 
                    at (0.1, -0.2) {$\mu_{3D}$ \\ $\Sigma_{3D}$};


                % ==========================================
                % 2. MAPPING / PROJECTION ARROW
                % ==========================================
                \draw[-Latex, line width=1.5mm, gray!60] (2.2, 0) -- (5.5, 0) 
                    node[midway, above, font=\large\bfseries, text=blue!80!black, align=center, yshift=2mm] {Projection \\ (Jacobian J)};


                % ==========================================
                % 3. 2D GAUSSIAN (IMAGE SPACE)
                % ==========================================
                \coordinate (mu2D) at (9, 0);

                % Scope for Image Plane with a gentle rotation
                \begin{scope}[shift={(mu2D)}, rotate=12] 
                    
                    % 16x16 Tile Grid Representation
                    \fill[blue!5, rounded corners=2pt] (-2.8, -2.8) rectangle (2.8, 2.8);
                    \draw[step=1.4cm, blue!20, thick] (-2.8, -2.8) grid (2.8, 2.8);
                    \draw[blue!60, very thick, rounded corners=2pt] (-2.8, -2.8) rectangle (2.8, 2.8);

                    % Label for Image Space (Rotation inherits from scope)
                    \node[blue!80!black, font=\Large\bfseries, anchor=south] at (0, 3) {2D Image Space};

                    % 2D Fuzzy Blob (Flattened via projection)
                    \foreach \i in {1,2,3,4,5,6} {
                        \pgfmathsetmacro{\opac}{0.85 - \i*0.13}
                        \fill[green!50!black, opacity=\opac, rotate around={35:(0,0)}]
                            (0,0) ellipse ({1.7 - \i*0.25} and {0.7 - \i*0.1});
                    }

                    % 2D Coordinate Stub representing Covariance (\Sigma_{2D})
                    % Arrows point strictly along the 2D principal axes
                    \draw[-Latex, very thick, red!80!black] (0,0) -- ++(35:1.6) node[right, font=\scriptsize] {}; % Major axis U
                    \draw[-Latex, very thick, red!80!black] (0,0) -- ++(125:0.6) node[above, font=\scriptsize] {}; % Minor axis V

                    % Center Point and Label
                    \node[circle, fill=black, inner sep=1.8pt] at (0,0) {};
                    \node[below right, font=\small\bfseries, align=left, fill=white, fill opacity=0.7, text opacity=1, inner sep=2pt, rounded corners] 
                        at (0.1, -0.2) {$\mu_{2D}$ \\ $\Sigma_{2D}$};
                        
                \end{scope}

            \end{tikzpicture}
        \end{figure}
    }
\end{frame}

\begingroup
    \addtocounter{framenumber}{-1}
    \againframe{forwardpass_pipeline}
\endgroup

\begin{frame}{Forward Pass: Phases 4 \& 5}
    \only<1-2>{
        \textbf{Phase 4: Tile Mapping (Preparing Radix Sort)}
        \begin{itemize}
            \justifying
            \item The target output image is divided into a grid of $16\times16$ pixel tiles.
            \item The algorithm calculates a 2D bounding box to map exactly which tiles a given 2D Gaussian covers.
            \item Assigns each 2D Gaussian to a tile. If overlapping, assigns to all affected tiles.
        \end{itemize}
        
        \vspace{1em}
        
        \uncover<2->{
            \textbf{Phase 5: Global Radix Sort}
            \begin{itemize}
                \justifying
                \item The GPU waits until all per-Gaussian computations are finished.
                \item A highly optimized Radix Sort is executed, organizing the 2D Gaussians for each tile strictly by depth (closest to farthest).
            \end{itemize}
        }
    }

    \only<3>{
        \begin{figure}
            \centering
            % Scales the entire graphic to perfectly fit the slide width
            \resizebox{0.6\linewidth}{!}{%
            \begin{tikzpicture}[every node/.style={transform shape}]
                
                % 1. BACKGROUND FILLS
                % Main 16x16 Tile Background
                \fill[blue!5] (0, 0) rectangle (8, -8);
                % Adjacent Tiles Backgrounds
                \fill[gray!10] (8, 0) rectangle (9, -8); % Right
                \fill[gray!10] (0, -8) rectangle (8, -9); % Bottom
                \fill[gray!10] (8, -8) rectangle (9, -9); % Bottom-Right
                
                % 2. PIXEL GRID (1 pixel = 0.5 units -> 16 pixels = 8 units)
                % Draws the 18x18 pixel grid
                \draw[step=0.5cm, blue!20, very thin] (0, -9) grid (9, 0);
                
                % 3. TILE BOUNDARIES
                % Main 16x16 Tile Border
                \draw[blue!80!black, very thick] (0, 0) rectangle (8, -8);
                % Adjacent Tile Separators
                \draw[gray!50, very thick, dashed] (8, -8) -- (8, -9);
                \draw[gray!50, very thick, dashed] (8, -8) -- (9, -8);
                
                % 4. LABELS FOR TILES
                \node[blue!80!black, font=\Large\bfseries, anchor=north west] at (0.2, -0.2) {16$\times$16 Pixel Tile};
                \node[gray!60!black, font=\small\bfseries, anchor=north west] at (0.2, -8.225) {Adjacent Tiles};

                
                % ==========================================
                % 5. 2D GAUSSIANS & BOUNDING BOXES
                % ==========================================

                % --- GAUSSIAN 0: Behind Gaussian 1 (Fully Inside) ---
                % Center: (4.2, -1.8), Radii: (1.4, 1.0), Rotation: -20 degrees
                \foreach \s in {1.0, 0.93, 0.87, 0.8, 0.73, 0.67, 0.6, 0.53, 0.47, 0.4, 0.33, 0.27, 0.2, 0.13, 0.07} {
                    \fill[cyan!60!blue, opacity=0.08, rotate around={-20:(4.2,-2.8)}] (4.2,-2.8) ellipse ({1.4*\s} and {1.0*\s});
                }
                \draw[cyan!60!blue, thick, rotate around={-20:(4.2,-2.8)}] (4.2,-2.8) ellipse (1.4 and 1.0);
                
                % --- GAUSSIAN 1: Fully Inside ---
                % Center: (3, -3), Radii: (1.8, 0.9), Rotation: 35 degrees
                \foreach \s in {1.0, 0.93, 0.87, 0.8, 0.73, 0.67, 0.6, 0.53, 0.47, 0.4, 0.33, 0.27, 0.2, 0.13, 0.07} {
                    \fill[green!50!black, opacity=0.08, rotate around={35:(3,-4)}] (3,-4) ellipse ({1.8*\s} and {0.9*\s});
                }
                \draw[green!50!black, thick, rotate around={35:(3,-4)}] (3,-4) ellipse (1.8 and 0.9);
                % Calculated tight bounding box
                \draw[red!80!black, thick, dashed] (1.44, -2.73) rectangle (4.56, -5.27);
                % Calculated tight bounding box
                \draw[red!80!black, thick, dashed] (2.84, -1.75) rectangle (5.56, -3.85);
                \node[red!80!black, font=\footnotesize\bfseries, anchor=north] at (2.6, -5.27) {Bounding Boxes};


                % --- GAUSSIAN 2: Overlapping Boundary ---
                % Center: (7.0, -6.0), Radii: (1.5, 0.8), Rotation: 45 degrees
                \foreach \s in {1.0, 0.93, 0.87, 0.8, 0.73, 0.67, 0.6, 0.53, 0.47, 0.4, 0.33, 0.27, 0.2, 0.13, 0.07} {
                    \fill[purple!60!black, opacity=0.08, rotate around={45:(7.0,-6.0)}] (7.0,-6.0) ellipse ({1.5*\s} and {0.8*\s});
                }
                \draw[purple!60!black, thick, rotate around={45:(7.0,-6.0)}] (7.0,-6.0) ellipse (1.5 and 0.8);
                % Calculated tight bounding box (bleeds past X=8)
                \draw[red!80!black, thick, dashed] (5.8, -4.8) rectangle (8.2, -7.2);
                
                % Annotation for the boundary overlap
                \draw[-Latex, thick, purple!80!black] (6.5, -7.6) -- (8.1, -7.3);
                \node[font=\footnotesize\bfseries, purple!80!black, align=center, anchor=north] at (6.2, -7.55) {Box overlaps boundary\\(assigned to both tiles)};


                % --- GAUSSIAN 3: Completely Outside ---
                % Center: (8.5, -2), Radii: (0.3, 0.6), Rotation: 15 degrees
                \foreach \s in {1.0, 0.93, 0.87, 0.8, 0.73, 0.67, 0.6, 0.53, 0.47, 0.4, 0.33, 0.27, 0.2, 0.13, 0.07} {
                    \fill[gray!70!black, opacity=0.06, rotate around={15:(8.5,-2)}] (8.5,-2) ellipse ({0.3*\s} and {0.6*\s});
                }
                \draw[gray!70!black, thick, rotate around={15:(8.5,-2)}] (8.5,-2) ellipse (0.3 and 0.6);
                % Calculated tight bounding box (strictly X > 8)
                \draw[red!80!black, thick, dashed] (8.17, -1.42) rectangle (8.83, -2.58);
                
                % Annotation for excluded Gaussian
                \draw[-Latex, thick, gray!80!black] (9.0, -0.74) -- (8.5, -1.3); 
                \node[font=\footnotesize\bfseries, gray!80!black, align=center, anchor=south] at (9.2, -0.9) {Excluded\\(outside tile)};
                
            \end{tikzpicture}%
            }
        \end{figure}
    }
\end{frame}

\begin{frame}{Forward Pass: Phase 6}
    % Overlay 1: Text only
    % \only<1>{
    \textbf{Phase 6: Tile-Based Rasterization (Splatting)}
    \vspace*{2mm}
    \begin{itemize}
        \justifying
        \item Operating in 2D image space, the GPU assigns one processing thread per pixel.
        \item \textbf{Alpha Compositing:} Pixels accumulate color by blending the sorted list of overlapping 2D Gaussians front-to-back:\\
        \vspace*{2mm}
        {\hspace{1.1cm}
        \resizebox{0.75\linewidth}{!}{
        $$
            \darkmathhighlight{
            C_{p} = \sum_{i=1}^{N} c_i \alpha_i T_i,
            \hspace{0.5cm} \text{with} \hspace{0.5cm}
            T_i = \prod_{j=1}^{i-1} (1 - \alpha_j)
            \hspace{0.5cm} \text{the transmittance}}
        $$
        }}
    \end{itemize}
        %\item \textbf{Early Stopping:} Threads stop calculating exactly when transmittance ($T_i$) drops near zero, ignoring hidden geometry for massive speed-ups.
    \pause
    \vspace*{1mm}
    \begin{figure}
        \centering
        \begin{tikzpicture}[scale=0.65, every node/.style={transform shape}]
            
            % 1. Mathematically derived blended color from alpha compositing
            % C_p = (0.4 * Red) + (0.5 * 0.6 * Green) + (0.8 * 0.3 * Blue)
            \definecolor{blendedColor}{rgb}{0.64, 0.54, 0.0}
            
            % 2. The Target Pixel on the View Plane
            \draw[thick, blue!60!black, fill=blendedColor] (-2.5, -1.2) rectangle (-0.5, 1.2);
            \node[font=\Large\bfseries, blue!80!black, align=center] at (-1.5, 1.9) {Target \\ Pixel $p$};
            \node[circle, fill=black, inner sep=1.5pt] (pixel) at (-1.5,0) {};
            
            % Labels inside the colored pixel
            \node[font=\small\bfseries, text=white] at (-1.5, 0.4) {Blended};
            \node[font=\small\bfseries, text=white] at (-1.5, -0.4) {Color};
            
            % 3. The Rendering Thread / Depth Ray
            \draw[-Latex, very thick, gray] (pixel) -- (16.5, 0) node[right, font=\Large\bfseries] {Depth};
            
            % 4. Sorted 2D Gaussians (Alpha Compositing)
            % Splat 1 (Target alpha 0.4 -> ~0.09 per layer)
            \foreach \s in {1, 0.8, 0.6, 0.4, 0.2} {
                \fill[red, opacity=0.09] (2.5, 0) ellipse ({0.6*\s} and {1.2*\s});
            }
            \draw[red!80!black, thick] (2.5, 0) ellipse (0.6 and 1.2);
            \node[circle, fill=black!80, inner sep=1.2pt] at (2.5, 0) {}; % Piercing dot
            \node[above, font=\large\bfseries] at (2.5, 1.4) {Gaussian 1};
            \node[below, font=\normalsize, align=center] at (2.5, -1.3) {Opacity \\ $\alpha_1=0.4$};
            \node[below, font=\normalsize\bfseries, text=blue!70!black] at (2.5, -2.3) {$T_1=1.00$};
            
            % Splat 2 (Target alpha 0.5 -> ~0.12 per layer)
            \foreach \s in {1, 0.8, 0.6, 0.4, 0.2} {
                \fill[green, opacity=0.12] (6, 0.3) ellipse ({0.7*\s} and {1.3*\s});
            }
            \draw[green!80!black, thick] (6, 0.3) ellipse (0.7 and 1.3);
            \node[circle, fill=black!80, inner sep=1.2pt] at (6, 0) {}; % Piercing dot
            \node[above, font=\large\bfseries] at (6, 1.8) {Gaussian 2};
            \node[below, font=\normalsize, align=center] at (6, -1.3) {Opacity \\ $\alpha_2=0.5$};
            \node[below, font=\normalsize\bfseries, text=blue!70!black] at (6, -2.3) {$T_2=0.60$};
            
            % Splat 3 (Target alpha 0.8 -> ~0.27 per layer)
            \foreach \s in {1, 0.8, 0.6, 0.4, 0.2} {
                \fill[yellow, opacity=0.27] (9.5, -0.2) ellipse ({0.5*\s} and {1.0*\s});
            }
            \draw[yellow!80!black, thick] (9.5, -0.2) ellipse (0.5 and 1.0);
            \node[circle, fill=black!80, inner sep=1.2pt] at (9.5, 0) {}; % Piercing dot
            \node[above, font=\large\bfseries] at (9.5, 1.2) {Gaussian 3};
            \node[below, font=\normalsize, align=center] at (9.5, -1.3) {Opacity \\ $\alpha_3=0.8$};
            \node[below, font=\normalsize\bfseries, text=blue!70!black] at (9.5, -2.3) {$T_3=0.30$};
            
            % 5. Early Stopping Barrier
            \draw[dashed, red, very thick] (11.2, -2.5) -- (11.2, 2.5);
            \node[red, font=\Large\bfseries, rotate=90] at (11.7, 0) {Early Stopping};
            \node[red, font=\normalsize\bfseries, rotate=90] at (12.2, -0.2) {($T_4 \approx 0$)};
            
            % 6. Culled / Ignored Geometry
            % Splat 4 (Ignored)
            \foreach \s in {1, 0.8, 0.6, 0.4, 0.2} {
                \fill[gray, opacity=0.07] (13.5, 0.4) ellipse ({0.6*\s} and {1.1*\s});
            }
            \draw[gray, dashed, thick] (13.5, 0.4) ellipse (0.6 and 1.1);
            \node[circle, fill=gray!80!black, inner sep=1.2pt] at (13.5, 0) {}; % Piercing dot
            \node[above, font=\large\bfseries, text=gray] at (13.5, 1.7) {Gaussian 4};
            
            % Splat 5 (Ignored)
            \foreach \s in {1, 0.8, 0.6, 0.4, 0.2} {
                \fill[gray, opacity=0.04] (15, -0.3) ellipse ({0.7*\s} and {0.9*\s});
            }
            \draw[gray, dashed, thick] (15, -0.3) ellipse (0.7 and 0.9);
            \node[circle, fill=gray!80!black, inner sep=1.2pt] at (15, 0) {}; % Piercing dot
            \node[above, font=\large\bfseries, text=gray] at (15.2, 0.8) {Gaussian 5};
            
            \node[text=gray, font=\Large\bfseries, align=center] at (14.25, -1.9) {Geometry\\Ignored};
            
        \end{tikzpicture}%
    \end{figure}
    % }

    \pdfcomment[opacity=0]{
        - mention GPU thread blocks \textCR
        - mention early stopping (the transmittance enables early stopping etc) \textCR
        \textCR
        - This completes the forward pass (continue to training cycle)
    }
    
    % Overlay 2: Graphic only (replaces the text on the same frame number)
    % \only<2>{
    %     \begin{figure}
    %         \centering
    %         % Scales the entire graphic to perfectly fit the slide width
    %         \hspace*{-5mm}
    %         \begin{tikzpicture}[scale=0.73, every node/.style={transform shape}]
                
    %             % 1. Mathematically derived blended color from alpha compositing
    %             % C_p = (0.4 * Red) + (0.5 * 0.6 * Green) + (0.8 * 0.3 * Blue)
    %             \definecolor{blendedColor}{rgb}{0.64, 0.54, 0.0}
                
    %             % 2. The Target Pixel on the View Plane
    %             \draw[thick, blue!60!black, fill=blendedColor] (-2.5, -1.2) rectangle (-0.5, 1.2);
    %             \node[font=\Large\bfseries, blue!80!black, align=center] at (-1.5, 1.9) {Target \\ Pixel $p$};
    %             \node[circle, fill=black, inner sep=1.5pt] (pixel) at (-1.5,0) {};
                
    %             % Labels inside the colored pixel
    %             \node[font=\small\bfseries, text=white] at (-1.5, 0.4) {Blended};
    %             \node[font=\small\bfseries, text=white] at (-1.5, -0.4) {Color};
                
    %             % 3. The Rendering Thread / Depth Ray
    %             \draw[-Latex, very thick, gray] (pixel) -- (16.5, 0) node[right, font=\Large\bfseries] {Depth};
                
    %             % 4. Sorted 2D Gaussians (Alpha Compositing)
    %             % Splat 1 (Target alpha 0.4 -> ~0.09 per layer)
    %             \foreach \s in {1, 0.8, 0.6, 0.4, 0.2} {
    %                 \fill[red, opacity=0.09] (2.5, 0) ellipse ({0.6*\s} and {1.2*\s});
    %             }
    %             \draw[red!80!black, thick] (2.5, 0) ellipse (0.6 and 1.2);
    %             \node[circle, fill=black!80, inner sep=1.2pt] at (2.5, 0) {}; % Piercing dot
    %             \node[above, font=\large\bfseries] at (2.5, 1.4) {Gaussian 1};
    %             \node[below, font=\normalsize, align=center] at (2.5, -1.3) {Opacity \\ $\alpha_1=0.4$};
    %             \node[below, font=\normalsize\bfseries, text=blue!70!black] at (2.5, -2.3) {$T_1=1.00$};
                
    %             % Splat 2 (Target alpha 0.5 -> ~0.12 per layer)
    %             \foreach \s in {1, 0.8, 0.6, 0.4, 0.2} {
    %                 \fill[green, opacity=0.12] (6, 0.3) ellipse ({0.7*\s} and {1.3*\s});
    %             }
    %             \draw[green!80!black, thick] (6, 0.3) ellipse (0.7 and 1.3);
    %             \node[circle, fill=black!80, inner sep=1.2pt] at (6, 0) {}; % Piercing dot
    %             \node[above, font=\large\bfseries] at (6, 1.8) {Gaussian 2};
    %             \node[below, font=\normalsize, align=center] at (6, -1.3) {Opacity \\ $\alpha_2=0.5$};
    %             \node[below, font=\normalsize\bfseries, text=blue!70!black] at (6, -2.3) {$T_2=0.60$};
                
    %             % Splat 3 (Target alpha 0.8 -> ~0.27 per layer)
    %             \foreach \s in {1, 0.8, 0.6, 0.4, 0.2} {
    %                 \fill[yellow, opacity=0.27] (9.5, -0.2) ellipse ({0.5*\s} and {1.0*\s});
    %             }
    %             \draw[yellow!80!black, thick] (9.5, -0.2) ellipse (0.5 and 1.0);
    %             \node[circle, fill=black!80, inner sep=1.2pt] at (9.5, 0) {}; % Piercing dot
    %             \node[above, font=\large\bfseries] at (9.5, 1.2) {Gaussian 3};
    %             \node[below, font=\normalsize, align=center] at (9.5, -1.3) {Opacity \\ $\alpha_3=0.8$};
    %             \node[below, font=\normalsize\bfseries, text=blue!70!black] at (9.5, -2.3) {$T_3=0.30$};
                
    %             % 5. Early Stopping Barrier
    %             \draw[dashed, red, very thick] (11.5, -2.5) -- (11.5, 2.5);
    %             \node[red, font=\Large\bfseries, rotate=90] at (12, 0) {Early Stopping ($T_4 = 0.06 \approx 0$)};
                
    %             % 6. Culled / Ignored Geometry
    %             % Splat 4 (Ignored)
    %             \foreach \s in {1, 0.8, 0.6, 0.4, 0.2} {
    %                 \fill[gray, opacity=0.07] (13.5, 0.4) ellipse ({0.6*\s} and {1.1*\s});
    %             }
    %             \draw[gray, dashed, thick] (13.5, 0.4) ellipse (0.6 and 1.1);
    %             \node[circle, fill=gray!80!black, inner sep=1.2pt] at (13.5, 0) {}; % Piercing dot
    %             \node[above, font=\large\bfseries, text=gray] at (13.5, 1.7) {Gaussian 4};
                
    %             % Splat 5 (Ignored)
    %             \foreach \s in {1, 0.8, 0.6, 0.4, 0.2} {
    %                 \fill[gray, opacity=0.04] (15, -0.3) ellipse ({0.7*\s} and {0.9*\s});
    %             }
    %             \draw[gray, dashed, thick] (15, -0.3) ellipse (0.7 and 0.9);
    %             \node[circle, fill=gray!80!black, inner sep=1.2pt] at (15, 0) {}; % Piercing dot
    %             \node[above, font=\large\bfseries, text=gray] at (15.2, 0.8) {Gaussian 5};
                
    %             \node[text=gray, font=\Large\bfseries, align=center] at (14.25, -1.9) {Geometry\\Ignored};
                
    %         \end{tikzpicture}%
    %     \end{figure}
    % }
\end{frame}

\begingroup
    \addtocounter{framenumber}{-1}
    \againframe{training_cycle}
\endgroup
